% Options for packages loaded elsewhere
\PassOptionsToPackage{unicode}{hyperref}
\PassOptionsToPackage{hyphens}{url}
\documentclass[
  ignorenonframetext,
]{beamer}
\newif\ifbibliography
\usepackage{pgfpages}
\setbeamertemplate{caption}[numbered]
\setbeamertemplate{caption label separator}{: }
\setbeamercolor{caption name}{fg=normal text.fg}
\beamertemplatenavigationsymbolsempty
% remove section numbering
\setbeamertemplate{part page}{
  \centering
  \begin{beamercolorbox}[sep=16pt,center]{part title}
    \usebeamerfont{part title}\insertpart\par
  \end{beamercolorbox}
}
\setbeamertemplate{section page}{
  \centering
  \begin{beamercolorbox}[sep=12pt,center]{section title}
    \usebeamerfont{section title}\insertsection\par
  \end{beamercolorbox}
}
\setbeamertemplate{subsection page}{
  \centering
  \begin{beamercolorbox}[sep=8pt,center]{subsection title}
    \usebeamerfont{subsection title}\insertsubsection\par
  \end{beamercolorbox}
}
% Prevent slide breaks in the middle of a paragraph
\widowpenalties 1 10000
\raggedbottom
\AtBeginPart{
  \frame{\partpage}
}
\AtBeginSection{
  \ifbibliography
  \else
    \frame{\sectionpage}
  \fi
}
\AtBeginSubsection{
  \frame{\subsectionpage}
}
\usepackage{iftex}
\ifPDFTeX
  \usepackage[T1]{fontenc}
  \usepackage[utf8]{inputenc}
  \usepackage{textcomp} % provide euro and other symbols
\else % if luatex or xetex
  \usepackage{unicode-math} % this also loads fontspec
  \defaultfontfeatures{Scale=MatchLowercase}
  \defaultfontfeatures[\rmfamily]{Ligatures=TeX,Scale=1}
\fi
\usepackage{lmodern}
\ifPDFTeX\else
  % xetex/luatex font selection
\fi
% Use upquote if available, for straight quotes in verbatim environments
\IfFileExists{upquote.sty}{\usepackage{upquote}}{}
\IfFileExists{microtype.sty}{% use microtype if available
  \usepackage[]{microtype}
  \UseMicrotypeSet[protrusion]{basicmath} % disable protrusion for tt fonts
}{}
\makeatletter
\@ifundefined{KOMAClassName}{% if non-KOMA class
  \IfFileExists{parskip.sty}{%
    \usepackage{parskip}
  }{% else
    \setlength{\parindent}{0pt}
    \setlength{\parskip}{6pt plus 2pt minus 1pt}}
}{% if KOMA class
  \KOMAoptions{parskip=half}}
\makeatother
\setlength{\emergencystretch}{3em} % prevent overfull lines
\providecommand{\tightlist}{%
  \setlength{\itemsep}{0pt}\setlength{\parskip}{0pt}}
\usepackage{bookmark}
\IfFileExists{xurl.sty}{\usepackage{xurl}}{} % add URL line breaks if available
\urlstyle{same}
\hypersetup{
  hidelinks,
  pdfcreator={LaTeX via pandoc}}

\author{\texorpdfstring{}{}}
\date{}

\begin{document}

\section{Introduction: Evidence Gaps in a Rapidly Expanding
Field}\label{introduction-evidence-gaps-in-a-rapidly-expanding-field}

\begin{frame}{The Free Energy Principle and Active Inference Framework}
\protect\phantomsection\label{the-free-energy-principle-and-active-inference-framework}
The Free Energy Principle (FEP), introduced by Karl Friston, proposes
that self-organizing systems maintain their structural and functional
integrity by minimizing variational free energy---an upper bound on
sensory surprise \citep{friston2006free, friston2010free}. Under this
principle, living systems are cast as approximate Bayesian inference
engines that build generative models of their environment and act to
reduce the discrepancy between predicted and observed states. Active
Inference (AIF) extends this picture from passive perception to
goal-directed behavior: agents select actions that bring about
observations consistent with their preferred states, unifying
perception, learning, and decision-making within a single variational
framework \citep{parr2022active, friston2017active}. Since its initial
formulation for sensorimotor control, AIF has been applied to
navigation, visual foraging, language comprehension, social cognition,
and multi-agent coordination.
\end{frame}

\begin{frame}{Challenges Posed by Rapid Literature Growth}
\protect\phantomsection\label{challenges-posed-by-rapid-literature-growth}
The active inference literature has expanded exponentially over the past
two decades, sustaining peak publication volumes into the late 2020s.
While early research concentrated almost exclusively on theoretical
neuroscience, the field has since diversified across biology (C5),
robotics (C2), computational psychiatry (C4), algorithm scaling (B), and
formal mathematics (A1). This rapid, multi-disciplinary growth creates
three interrelated challenges. First, tracking which core theoretical
claims---such as FEP universality or the physical realism of Markov
blankets---are deeply supported, contested, or merely assumed becomes
intractable. Second, because the relationship between mathematical
formalisms and empirical evidence remains frequently implicit,
systematic evidence synthesis demands prohibitive manual labor. Third,
new entrants must navigate a literature heavily weighted toward broad
qualitative philosophy (A2), interspersed with rapidly accelerating,
highly specialized applied pockets.

Traditional narrative reviews attempt to address these challenges but
are inherently static, subjective, and quickly outdated. Systematic
reviews from evidence-based medicine offer rigorous aggregation but are
structurally customized for clinical trial data with homogeneous outcome
measures, rendering them ill-suited for the heterogeneous ontological
and computational claims endemic to this theoretical literature. The
expansion of predictive processing
\citep{clark2013whatever, hohwy2013predictive} and the emergence of
formal parameterizations like Bayesian mechanics
\citep{sakthivadivel2023bayesian} further broaden the scope of
assertions that any comprehensive meta-analysis must reconcile.
\end{frame}

\begin{frame}{Related Work and Prior Meta-Analyses}
\protect\phantomsection\label{related-work-and-prior-meta-analyses}
Several prior efforts have surveyed aspects of the Active Inference
landscape. Sajid et al.~\citep{sajid2021active} compare active inference
with alternative decision-making frameworks; Da Costa et
al.~\citep{dacosta2020active} synthesize the discrete-state-space
formulation; Lanillos et al.~\citep{lanillos2021active} survey robotics
applications; Smith et al.~\citep{smith2021computational} provide a
tutorial bridging theory and empirical data; and Millidge et
al.~\citep{millidge2021understanding} examine information-theoretic
foundations of exploration behavior. Ramstead et
al.~\citep{ramstead2018answering} extend the FEP to questions of
biological self-organization, while Pezzulo et
al.~\citep{pezzulo2015active} connect active inference to homeostatic
regulation.

Closest to our work, Knight, Cordes, and Friedman \citep{knight2022fep}
conducted a systematic literature analysis of publications using the
terms ``Free Energy Principle'' or ``Active Inference,'' with an
emphasis on works by Karl J. Friston. Their analysis---maintained by the
Active Inference Institute---combined manual annotation of structural,
visual, and mathematical features with automated analyses using the
Active Inference Ontology at the scale of thousands of citations and
hundreds of annotated papers. That study identified six development
directions---including broader scope, richer annotation, and
transferable approaches---and represents an important precursor to
automated meta-analysis of this field.

These works are primarily narrative reviews: they synthesize qualitative
findings but do not strictly quantify the balance of evidence across the
field's central claims. The systematic analysis of Knight et
al.~\citep{knight2022fep} pioneered quantitative literature analysis for
this field using manual annotation and ontology-based automated
analysis. Our framework advances this line of work by (1) fully
automating assertion extraction via LLM-based hypothesis scoring, (2)
constructing a structured, RDF-compatible knowledge graph scored by
citation-weighted evidence, and (3) tracking how evidence for core
claims evolves over time through temporal trend analysis.
\end{frame}

\begin{frame}{Synergizing Knowledge Graphs and LLMs}
\protect\phantomsection\label{synergizing-knowledge-graphs-and-llms}
Recent systematic literature initiatives underscore a powerful
reciprocal synergy between Large Language Models (LLMs) and Knowledge
Graphs: LLMs parse unstructured text to rapidly extract semantic claims,
efficiently populating the structured, queryable architecture of the
graph \citep{quevedo2025combining, li2024unifying}. We adopt the
\emph{nanopublication} \citep{groth2010anatomy}---a minimal,
machine-readable unit of scientific evidence comprising a core assertion
bound to explicit provenance metadata---as the fundamental serialization
format for this extracted knowledge.
\end{frame}

\begin{frame}{This Study: Approach and Overview}
\protect\phantomsection\label{this-study-approach-and-overview}
This paper presents a computational meta-analysis of the Active
Inference literature (\(N = {{CORPUS_SIZE}}\)). Rather than relying
exclusively on bibliometric metadata or slow manual coding, we deploy a
Large Language Model (LLM) to ``read'' each paper's abstract and assess
its relationship to eight core hypotheses within the FEP paradigm. We
serialize these assessments as nanopublications---each encoding an
assertion (``Paper X supports Hypothesis Y'') coupled with the LLM's
natural-language reasoning and confidence score. The resulting knowledge
graph aggregates these nanopublications and links them to paper
metadata, citation networks, subfield classifications, and hypothesis
definitions. A citation-weighted scoring formula quantifies the net
evidence for or against each hypothesis, producing scores in \([-1, 1]\)
that reflect both the direction and strength of published evidence.
\end{frame}

\begin{frame}{Research Questions}
\protect\phantomsection\label{research-questions}
This meta-analysis addresses four primary research questions:

\begin{enumerate}
\tightlist
\item
  \textbf{RQ1 (Field Structure):} What is the disciplinary structure and
  growth trajectory of the Active Inference literature, and how are
  papers distributed across the three domains---Core Theory (A), Tools
  \& Translation (B), and Application Domains (C)?
\item
  \textbf{RQ2 (Growth Dynamics):} What are the temporal growth dynamics
  of the field, and which subfields are experiencing the most rapid
  expansion?
\item
  \textbf{RQ3 (Hypothesis Evidence):} What is the current balance of
  evidence for and against the eight standard hypotheses, and how has
  this balance evolved over time? (See hypothesis dashboard and
  assertion figures in §6.)
\item
  \textbf{RQ4 (Tooling Readiness):} What is the state of software
  tooling and infrastructure for Active Inference research, and what
  gaps remain?
\end{enumerate}
\end{frame}

\begin{frame}{Scope and Delimitations}
\protect\phantomsection\label{scope-and-delimitations}
This study focuses on the English-language peer-reviewed and preprint
literature retrievable from arXiv, Semantic Scholar, and OpenAlex. We do
not include book chapters or monographs not indexed by these sources,
software documentation, or non-English publications. Domain
classification uses keyword matching rather than expert annotation---a
deliberate trade-off favoring reproducibility over precision, whose
consequences we quantify in Section 7. Hypothesis scoring relies on
LLM-extracted assertions; the fidelity and limitations of this approach
are examined in Section 3. The hypothesis definitions and domain
taxonomy are informed by, but not identical to, the Active Inference
Ontology used by Knight et al.~\citep{knight2022fep}; future alignment
would enable direct comparison with that earlier analysis.
\end{frame}

\begin{frame}{Principal Contributions}
\protect\phantomsection\label{principal-contributions}
This work makes five contributions:

\begin{enumerate}
\item
  \textbf{A multi-source retrieval and deduplication pipeline} for
  Active Inference literature, using a canonical identifier hierarchy
  across three academic databases.
\item
  \textbf{A nanopublication-based knowledge graph schema} encoding
  directed, confidence-scored assertions about eight core hypotheses
  with full provenance tracking.
\item
  \textbf{A quantitative field overview} characterizing the growth,
  domain distribution (A/B/C taxonomy), citation topology, and latent
  topic structure of the Active Inference literature.
\item
  \textbf{An LLM-based hypothesis scoring dashboard} that produces
  differentiated evidence profiles with temporal trend visualization.
\item
  \textbf{A tooling assessment} of the software ecosystem supporting
  Active Inference research, including the implemented extraction
  pipeline, existing software (pymdp, SPM, RxInfer.jl), and knowledge
  graph infrastructure.
\end{enumerate}

The remainder of this paper is organized as follows. Section 2 surveys
the tooling landscape (RQ4) and extraction pipeline. Section 3 documents
the LLM-based extraction pipeline in detail. Section 4 provides the
technical appendix (notation, abbreviations, hypothesis definitions).
Section 5 describes the methodology. Section 6 presents the hypothesis
evidence landscape (RQ3). Section 7 presents the field overview with
domain-level analysis (RQ1, RQ2), supplemented by detailed domain
analyses (§7a), text analytics (§7b), and citation network topology
(§7c). Section 8 concludes with limitations and future directions;
Section 8a provides community recommendations and open questions.
\end{frame}

\end{document}
